Let $p_r$ denote the $r$th prime, and write $g_r:=p_{r+1}-p_r$ for the prime gap.
For a positive integer $n$, let $P^+(n)$ be the largest prime factor of $n$.
The property in the problem is:
\[
\exists\ \text{at least two integers }n\ \text{with }p_r<n<p_{r+1}\ \text{and }P^+(n)<g_r.
\]
The question is whether this happens for infinitely many $r$.

Edge cases: if $g_r=2$ then there are no integers strictly between $p_r$ and $p_{r+1}$; if $g_r=3$ there is exactly one integer in between, so the property cannot hold. Thus any $r$ witnessing the property must have $g_r\ge 4$.
